\scanpage{22}
\mspage{- 14 -}

n\textsuperscript{o} 391

\[
(t,\lambda)\handwritten{\longrightarrow}({}^{v}u(t)+\lambda(u(a)-a'),\lambda).\tag*{(9)}
\]

Il est clair que l'application $u\handwritten{\longrightarrow}{}^{h}u$ est \uline{injective}.

On appelle \uline{pseudo-réflexion affine} une application affine $s$ de $E$ dans lui-même qui laisse fixes les points d'un \uline{hyperplan \unsure{affine|af'ine}} $H$. Si on prend pour origine dans $E$ un point $a\in H$, $s$ devient une \uline{pseudo-réflexion} au sens du n\textsuperscript{o} 2 ; donc $^{v}s$ est une pseudo-réflexion laissant invariants les vecteurs de l'hyperplan $^{v}H$, \uline{direction} de $H$ ; et la formule (9) montre que $^{h}s$ est une \uline{pseudo-réflexion} dans l'espace vectoriel $^{h}E$, laissant invariants les points de l'hyperplan $^{h}H$ \uline{engendré par} $H$.

On dit qu'une pseudo-réflexion affine $s$ est une \uline{réflexion affine} si elle est involutive et distincte de \deleted{l'}identité ; il revient au même de dire que la pseudo-réflexion $^{v}s$ ou $^{h}s$ est une \uline{réflexion} (n\textsuperscript{o} 2), comme il résulte aussitôt de ce qui précède.

Lorsque $K$ est de caractéristique $\neq2$ et que $T$ est muni d'une forme bilinéaire symétrique non dégénérée $\Phi$, les réflexions de $E$ qui sont des \uline{déplacements} pour la forme métrique $\Phi$ (\uline{Alg}., chap. IX, §6, n\textsuperscript{o} 6) sont les \uline{symétries orthogonales par rapport aux hyperplans affines} (non isotropes) de $E$ (n\textsuperscript{o} 2, prop. 4).

\margin{\handwritten{\unsure{/}}}

\margin{\handwritten{\unsure{!}}}

PROPOSITION 9.- \uline{Soient} $H$, $H'$ \uline{deux hyperplans dans} $E$ \uline{parallèles et non isotropes} (\uline{pour} $\Phi$) ; \uline{il existe un unique vecteur} $t\in T$ \uline{orthogonal à} $H$ \uline{et tel que} $H'=t+H$. \uline{Si} $s$, $s'$ \uline{sont les symétries orthogonales par rapport à} $H$ \uline{et} $H'$ \uline{respectivement}, $s's$ \uline{est la translation} $2t$ \uline{dans} $E$.

En prenant dans $E$ une origine contenue dans $H$, on peut se borner au cas où $E=T$. L'existence et l'unicité de $t$ sont immédiates. En prenant une nouvelle origine au point $t$, on voit que

\[
s'(x)-t=s(x-t)=s(x)+t
\]

% archive sheet number: 21
